\section{Discussion}
\label{sec:discussion}

\subsection{Where it worked, and why}

LOO temperature scaling alone is a real, if not universal, calibration improvement. The
mechanism is direct: it always smooths the model's confidence toward a single label-free
accuracy estimate computed from the cache. Where TDA's cache was badly overconfident relative
to its own accuracy -- DTD (17.41\% $\rightarrow$ 4.68\% ECE) and Aircraft (15.56\%
$\rightarrow$ 4.21\%, same-code-path \texttt{none}-vs-\texttt{temp} comparison) chief among
them -- pulling confidence down toward the estimate is exactly the correction needed, and the
improvement is large.

\subsection{Where it did not, and why}

The same mechanism explains the Pets reversal. On Pets, TDA's cache was already close to
perfectly calibrated (2.52\% ECE for the headline \texttt{tda} run, 2.48\% for the ablation's
\texttt{none} baseline) -- the model's stated confidence already matched its
accuracy. LOO temperature scaling has no notion of ``already well calibrated''; it always
smooths toward the LOO estimate regardless of how good the starting point was, so with no
overconfidence to correct, it manufactures underconfidence instead (ECE rises to 15.54\% under
\texttt{temp}, or as high as 43.98\% under the full bundle). A method that only ever pushes
confidence toward one estimate, with no mechanism to detect when the starting point does not
need correcting, will help wherever that estimate happens to be closer to the truth than the
model's own confidence, and hurt wherever the estimate is not.

Margin admission and probabilistic fusion fail for related but distinct reasons. Margin
admission's single global threshold (\texttt{margin\_threshold} $= 1.0$) is not portable:
measured admission rates (Section~\ref{sec:admission-rates}) run from 15.9\% on Aircraft to
81.9\% on Pets, and they track zero-shot CLIP's accuracy rather than class count. On Aircraft the cache
mostly starves; on Pets the gate is nearly a no-op. The gate is therefore most selective exactly
where the cache had the most to contribute. Neither
extreme is what a per-dataset-tuned gate would look like, and, as Section~\ref{sec:results} shows, its
selectivity costs accuracy outright (56.60\% $\rightarrow$ 53.67\% under margin admission
alone).

Probabilistic fusion, meanwhile, is the single worst component in the ablation (mean ECE
24.54\%) -- but not for the reason it first appears. \texttt{compute\_cache\_logits}
(inherited unmodified from upstream TDA) returns, for each class, a sum of
$\exp(-\beta(1 - \mathrm{affinity}))$ terms over whatever cached items belong to that class --
an unnormalized, comparatively small-magnitude quantity, since affinity is a cosine similarity
in $[-1, 1]$ and only the handful of classes actually represented in a small cache get a
nonzero score at all. $\sigma(z_{\text{cache}})$ is therefore close to uniform relative to
$\sigma(z_{\text{clip}})$, which is sharply peaked because CLIP's own logits are already scaled
by its learned logit scale ($\approx 100$). Averaging a near-uniform distribution into a
near-one-hot one at $w = 0.5$ (Eq.~\eqref{eq:fusion}) caps top-1 probability well below what CLIP
alone would report, \emph{regardless of how much evidence actually supports the prediction} -- fusion
does not need both sources to agree on a wrong answer to produce this effect; it produces the effect on
every sample, right or wrong alike. This single fact accounts for everything the reliability
diagrams show (Section~\ref{sec:results}): no \texttt{cal\_tda} reliability record on any
dataset reaches even 0.70 mean confidence in its highest-occupied bin, mean confidence sits
8.7--44 points below accuracy on every dataset, and the size of that gap tracks accuracy itself:
it is largest on the highest-accuracy dataset (Pets, 88.74\% accuracy capped near 44.77\%
confidence) and smallest on the lowest-accuracy one (Aircraft, 23.97\% accuracy capped near
15.32\%). That is what a fixed confidence ceiling does -- it costs a high-accuracy dataset far
more calibration error than a low-accuracy one. Fusion's ECE is therefore not overconfidence from agreeing wrong
answers; it is underconfidence manufactured by capping every prediction's confidence near one
fixed point, independent of correctness.

This also explains Aircraft's apparent improvement in Table~\ref{tab:main} without invoking any
real correction. The \texttt{none} baseline on Aircraft is overconfident to begin with (mean
confidence 39.77\% against 24.21\% accuracy, ECE 15.56\%) -- ordinary TDA-style overconfidence.
Adding fusion does not fix that overconfidence; it overshoots straight past it, capping
confidence at 15.2--15.4\% against 23.97--24.42\% accuracy (ECE 8.65--9.20\%). The resulting
ECE is smaller purely because Aircraft's true accuracy (23.97\%) happens to sit close to
fusion's confidence ceiling on this particular dataset -- an accident of
where the ceiling landed relative to Aircraft's accuracy, not evidence that fusion diagnosed and
corrected \texttt{none}'s overconfidence. On every other dataset, where accuracy is higher, the
same ceiling produces a much larger gap in the other direction, which is why fusion is the worst
component overall despite this one favourable-looking case.

\subsection{The margin$\times$temperature interaction, and its limits}

An earlier single-dataset (DTD) observation suggested a clean interaction, and the five-dataset
runs bear it out: once margin
admission and fusion are both active, the LOO temperature search pins at its grid boundary
$T = 1.0$ on 98--100\% of samples on every dataset, meaning the search wants to sharpen further
than the smoothing-only grid allows and Contribution 3 becomes effectively inert on top of Contributions 1
and 2 -- consistent with the ECE difference between \texttt{margin\_fusion} and \texttt{all}
being tiny everywhere (0.00--0.36 points). Testing that story against temperature on top of
margin \emph{alone} (without fusion in between) shows it does not generalize cleanly: pinning
drops to 75--95\% on four datasets, with a correspondingly larger ECE effect (DTD improves by
2.81 points, Flowers102 by 2.25), and \textbf{Pets breaks the pattern outright} -- its
\texttt{margin\_temp} run pins only 0.2\% of samples at the boundary, the opposite of inert, and
it is exactly on Pets, already well calibrated under margin alone (2.53\% ECE), that adding
temperature does the most damage of any margin-plus-temperature run (2.53\% $\rightarrow$
14.79\%; the larger temperature-induced regression in the ablation is Pets without margin,
Table~\ref{tab:temp-per-dataset}: 2.48\% $\rightarrow$ 15.54\%). The
single-dataset story pointed to one mechanism; five datasets show it holds strongly with fusion
present, holds more weakly for margin alone, and inverts on the one dataset a single-dataset
study would have missed. We report this as a caution about generalizing a mechanism from one
dataset, not as a discarded finding -- the underlying explanation (temperature is inert exactly
when the search is already pinned against its floor) still holds; what changes across datasets
is how often that precondition is met.

\subsection{Accuracy/calibration trade-off, stated plainly}

None of the three contributions, alone or combined, is a free lunch. Contribution 3 is the one
change in this project that costs nothing in accuracy (it structurally cannot, being a
monotonic rescaling) and buys a real mean calibration improvement -- but even that
improvement is not uniform, and reverses on the one dataset that was already well calibrated.
Contributions 1 and 2 cost accuracy, calibration, or both, on average, and the one apparent
exception -- fusion's lower ECE on Aircraft -- is not a benefit either, once its mechanism is
understood: it is the same confidence-capping failure landing, by the accident of that one
dataset's low accuracy, closer to correct than the overconfidence it replaced. The honest summary of this project's
central finding is that a training-free, label-free calibration fix is achievable, but it is one
targeted change (temperature scaling on a label-free accuracy estimate), not the three-part
bundle we set out to build.
